\section*{Hoofdstuk \ref{ch:g6:propdata} --- Evenredigheid en gegevens}

\begin{solution}{exo:g6:propdata:1}
Appels: \omterm{def:g3:division:remainder}{quotiënten} $\frac52 = 2.5$, $\frac{7.5}{3} = 2.5$,
$\frac{12.5}{5} = 2.5$ --- alle gelijk: \omterm{def:g6:propdata:def}{evenredig} (prijs $2.50$ per kg).

Leeftijd en lengte: $\frac{85}{2} = 42.5$ maar $\frac{103}{4} = 25.75$ --- niet
gelijk: niet \omterm{def:g6:propdata:def}{evenredig}.
\end{solution}

\begin{solution}{exo:g6:propdata:2}
Eén croissant: $4.80 \div 4 = 1.20$ euro. Zeven: $7 \times 1.20 = 8.40$ euro.
\end{solution}

\begin{solution}{exo:g6:propdata:3}
De constante: $18 \div 10 = 1.8$ euro per liter. Dan kost $25$ L
$25 \times 1.8 = 45$ euro; voor $72$ euro krijg je $72 \div 1.8 = 40$ L; $60$ L
kost $60 \times 1.8 = 108$ euro.
\end{solution}

\begin{solution}{exo:g6:propdata:4}
Voor $250$ km: $2.5$ keer de brandstof van $100$ km, dus $2.5 \times 6 = 15$ L.
Voor $350$ km: $3.5 \times 6 = 21$ L. Met $27$ L: $27 \div 6 = 4.5$ honderden
km, dus $450$ km.
\end{solution}

\begin{solution}{exo:g6:propdata:5}
$50\,\%$ van $84$: de \omterm{def:g3:division:half}{helft}, $42$. \quad
$25\,\%$ van $200$: een \omterm{def:g3:division:half}{kwart}, $50$. \quad
$10\,\%$ van $63$: een tiende, $6.3$. \quad
$30\,\%$ van $70$: $3 \times 7 = 21$.
\end{solution}

\begin{solution}{exo:g6:propdata:6}
$60\,\%$ van $450$: $\frac{60}{100} \times 450 = 270$ leerlingen eten in de
kantine; $450 - 270 = 180$ niet.
\end{solution}

\begin{solution}{exo:g6:propdata:7}
Minder dan $10$ boeken: dinsdag ($8$) en donderdag ($6$). Woensdag min donderdag:
$17 - 6 = 11$ boeken meer.
\end{solution}

\begin{solution}{exo:g6:propdata:8}
Een staafdiagram met vijf staven van hoogte $14$, $16$, $13$, $17$ en $20$ (een
streepje per $2$ of per $5$ graden werkt goed). Warmste dag: vrijdag ($20$
graden).
\end{solution}

\begin{solution}{exo:g6:propdata:9}
Voor $10$ personen vermenigvuldig je met $\frac{10}{6} = \frac53$: bloem
$450 \times \frac{10}{6} = 750$ g. Eieren: $3 \times \frac{10}{6} = 5$ ---
gelukkig een heel getal. (Was het dat niet, dan rond je naar \emph{boven} af om
genoeg te hebben.)
\end{solution}

\begin{solution}{exo:g6:propdata:10}
\emph{1.} In $2$ u: $48$ km. In een half uur: $12$ km. In $1$ u $30$:
$24 + 12 = 36$ km.

\emph{2.} $60 \div 24 = 2.5$ uur, dus $2$ u $30$ min.
\end{solution}

\begin{solution}{exo:g6:propdata:11}
Bal: korting $20\,\%$ van $15 = 3$ euro, nieuwe prijs $12$ euro. Racket: korting
$8$ euro, nieuwe prijs $32$ euro. De nieuwe prijs is in elk geval $80\,\%$ van de
oude: \omterm{def:g6:propdata:def}{evenredig}, met constante $0.8$.
\end{solution}

\begin{solution}{exo:g6:propdata:12}
Na $t$ uur heeft de dikke kaars $\frac t6$ van haar hoogte opgebrand: ze staat op
$1 - \frac t6$ van de beginhoogte; de dunne op $1 - \frac t4$. We willen
\[
1 - \frac t4 = \frac12 \left(1 - \frac t6\right)
\quad\text{dus}\quad
1 - \frac t4 = \frac12 - \frac t{12}.
\]
Dan is $\frac12 = \frac t4 - \frac t{12} = \frac{3t - t}{12} = \frac{t}{6}$, dus
$t = 3$ uur. Controle: na $3$ u staat de dikke kaars op $1 - \frac36 = \frac12$
en de dunne op $1 - \frac34 = \frac14$ --- inderdaad de \omterm{def:g3:division:half}{helft} daarvan.
\end{solution}

\begin{solution}{pb:g6:propdata:1}
\textbf{1.} $100\,000$ cm $= 1\,000$ m $= 1$ km: één centimeter op de kaart is
één kilometer op de grond.

\textbf{2.} $7.5$ cm $\rightarrow 7.5$ km tussen de dorpen; $12$ cm
$\rightarrow 12$ km meeroever.

\textbf{3.} $20$ km $\rightarrow 20$ cm op de kaart.

\textbf{4.} $5$ km $= 500\,000$ cm, dus is de schaal $1:500\,000$. De wandelkaart
($1:100\,000$) is de meest gedetailleerde: die gebruikt $5$ cm papier waar de
andere er maar $1$ cm aan besteedt.

\textbf{5.}
\begin{center}
\renewcommand{\arraystretch}{1.2}
\begin{tabular}{l|ccc}
kaart (cm) & $1$ & $2$ & $7.5$ \\
\hline
echt (km) & $1$ & $2$ & $7.5$ \\
\end{tabular}
\end{center}
Echte afstand $=$ kaartafstand $\times 1$ (in deze eenheden): altijd
\emph{dezelfde} vermenigvuldiger, en dat is precies de definitie van
\omterm{def:g6:propdata:def}{evenredigheid} (\cref{def:g6:propdata:def}). De constante is hier $1$ kilometer
per centimeter.

\textbf{6.} $4$ m $= 400$ cm en $400 \div 50 = 8$; $3$ m $= 300$ cm en
$300 \div 50 = 6$: de plattegrond toont een rechthoek van $8$ cm $\times$ $6$
cm.

\textbf{7.} Bed: $200 \div 50 = 4$ cm bij $90 \div 50 = 1.8$ cm. Deuropening:
$80 \div 50 = 1.6$ cm.

\textbf{8.} Echte \omterm{def:g6:measure:area}{oppervlakte}: $4 \times 3 = 12$ m$^2$. \omterm{def:g6:measure:area}{Oppervlakte} van de
plattegrond: $8 \times 6 = 48$ cm$^2$. Omgezet:
$12$ m$^2 = 12 \times 10\,000 = 120\,000$ cm$^2$, en
\[
120\,000 \div 48 = 2\,500 .
\]
Het \omterm{def:g3:division:remainder}{quotiënt} is niet $50$ maar $2\,500 = 50 \times 50$.

\textbf{9.} Eén cm$^2$ papier staat voor een echt vierkant van $50$ cm bij $50$
cm, en dat bevat $50 \times 50 = 2\,500$ cm$^2$. Worden lengtes dus door $50$
gedeeld, dan worden \omterm{def:g6:measure:area}{oppervlaktes} door $50 \times 50 = 2\,500$ gedeeld:
\omterm{def:g6:measure:area}{oppervlaktes} volgen het \emph{kwadraat} van de schaal.

\textbf{10.} $6 \times 2\,500 = 15\,000$ cm$^2$, en
$15\,000 \div 10\,000 = 1.5$ m$^2$.

\textbf{11.} Het plaatje suggereert dat er op zondag \emph{twee keer} zoveel
verkocht is (een staaf van dubbele hoogte). De waarheid: de stijging is $5$ euro
op $105$, en $5 \div 105 \approx 0.048$: ongeveer $5\,\%$ meer, niet $100\,\%$
meer. De eigenaar negeerde de waarschuwing om te controleren of de as bij $0$
begint (\cref{met:g6:propdata:read}).

\textbf{12.} Met de as vanaf $0$ stijgen de staven tot $105$ en $110$ kleine
eenheden: twee staven van bijna dezelfde hoogte, het eerlijke plaatje van een
\omterm{ex:g1:subtraction:difference}{verschil} van $5\,\%$.

\textbf{13.} Omhoog: de groei is $20$ postzegels vanaf een start van $40$:
$\frac{20}{40} = 50\,\%$. Omlaag: de daling is $20$ postzegels vanaf een start
van $60$: $\frac{20}{60} = \frac13 \approx 33\,\%$. Dezelfde $20$ postzegels,
verschillende beginwaarden --- een percentage is altijd een \omterm{def:g6:fractions:def}{breuk} \emph{van
iets}, en dat ``iets'' is veranderd.

\textbf{14.} Na $20\,\%$ korting: $100 - 20 = 80$ euro. De extra $10\,\%$ geldt
voor $80$: korting $8$ euro, eindprijs $72$ euro. Totale korting: $28$ euro op
$100$, dus $28\,\%$ --- geen $30\,\%$. De tweede korting werkte op de al
verlaagde prijs, dus zijn haar euro's kleiner: percentages van verschillende
hoeveelheden tel je niet op.

\textbf{15.} Op die kaart staat $1$ cm voor $5$ km, dus staat $1$ cm$^2$ voor
$5 \times 5 = 25$ km$^2$ (de regel van vraag 9). Het bos:
$3 \times 25 = 75$ km$^2$. Lengtes volgen de schaal; \omterm{def:g6:measure:area}{oppervlaktes} haar
\emph{kwadraat}.
\end{solution}
